 \chapter{Thermal Diffusion: Numerical Implementation}

\section{3D Heat Diffusion Equation}

The differential form of the heat diffusion equation for an isotropic medium is:
\begin{equation}
    \rho c_p \frac{\partial T}{\partial t} = \frac{\partial}{\partial x}\left( k \frac{\partial T}{\partial x} \right) + \frac{\partial}{\partial y}\left( k \frac{\partial T}{\partial y} \right) + \frac{\partial}{\partial z}\left( k \frac{\partial T}{\partial z} \right) + \dot{q}
\end{equation}

We integrate the heat diffusion equation over the control volume $\Delta V = \Delta x \Delta y \Delta z$ centered at node $P$ (see Figure \ref{fig:3d_cv}):

\begin{equation}
    \int_{\Delta t} \iiint_V \rho c_p \frac{\partial T}{\partial t} dV dt = \int_{\Delta t} \oint_A (k \nabla T) \cdot d\mathbf{A} dt + \int_{\Delta t} \iiint_V \dot{q} dV dt
\end{equation}

Applying a fully implicit time-stepping scheme and evaluating face fluxes via central differencing, we arrive at the general discrete algebraic form:
\begin{equation}
    a_P T_P = a_W T_W + a_E T_E + a_S T_S + a_N T_N + a_B T_B + a_T T_T + b
\end{equation}

Where the coefficients are defined to maintain physical consistency across 3D space:
\begin{itemize}
    \item $a_E = \frac{k_e \Delta y \Delta z}{\delta x_{PE}}, \quad a_W = \frac{k_w \Delta y \Delta z}{\delta x_{WP}}$
    \item $a_N = \frac{k_n \Delta x \Delta z}{\delta y_{PN}}, \quad a_S = \frac{k_s \Delta x \Delta z}{\delta y_{SP}}$
    \item $a_T = \frac{k_t \Delta x \Delta y}{\delta z_{PT}}, \quad a_B = \frac{k_b \Delta x \Delta y}{\delta z_{BP}}$
    \item $a_P^0 = \frac{\rho c_p \Delta V}{\Delta t}$ (Temporal inertia)[cite: 1]
    \item $a_P = \sum a_{nb} + a_P^0 - S_P \Delta V$[cite: 1]
    \item $b = S_u \Delta V + a_P^0 T_P^0$[cite: 1]
\end{itemize}

\section{Simplification: 1D Steady-State Conduction}
To verify the TDMA kernel today, we collapse the 3D general form under two assumptions:
\begin{enumerate}
    \item \textbf{Steady-State:} $\frac{\partial T}{\partial t} \rightarrow 0$, therefore $a_P^0 = 0$[cite: 1].
    \item \textbf{1D Heat Flow:} $\Delta y, \Delta z \rightarrow 1$ and $a_N, a_S, a_T, a_B = 0$[cite: 1].
\end{enumerate}

The system collapses to a tridiagonal form suitable for the TDMA solver:
\begin{equation}
    a_P T_P = a_W T_W + a_E T_E + b
\end{equation}

In our code implementation (\texttt{solver\_1d.f90}), this maps to the matrix indices $i$:
\begin{equation}
    \underbrace{-a_W}_{a_i} T_{i-1} + \underbrace{a_P}_{b_i} T_i + \underbrace{-a_E}_{c_i} T_{i+1} = \underbrace{b}_{d_i}
\end{equation}


\section{Extension to 2D: Point-Iterative vs. Line-Implicit Methods}

Transitioning from 1D to 2D steady-state conduction introduces significant computational bottlenecks if treated purely explicitly. For a 2D uniform grid ($\Delta x = \Delta y$), the continuous Laplace equation discretizes to a 5-point stencil. We evaluate two distinct numerical architectures for solving this system, prioritizing cache coherency and convergence acceleration.

\subsection{The Point-Iterative Baseline (Jacobi)}



The explicit Jacobi method updates each node independently using the localized average of its orthogonal neighbors:

\begin{equation}
    T_{i,j}^{new} = \frac{a_E T_{i+1,j} + a_W T_{i-1,j} + a_N T_{i,j+1} + a_S T_{i,j-1} + b}{a_P}
\end{equation}

For a uniform, homogeneous grid (where $\Delta x = \Delta y$ and conductivity $k$ is constant), the coefficients collapse to $a_{nb} = 1$ and $a_P = 4$, reducing the system to the localized average of its orthogonal neighbors:

\begin{equation}
    T_{i,j}^{new} = \frac{1}{4} \left( T_{i+1,j} + T_{i-1,j} + T_{i,j+1} + T_{i,j-1} \right)
\end{equation}

While mathematically simple and "embarrassingly parallel" (ideal for naive GPU porting), 
this Point-Iterative approach is highly inefficient. It possesses a spectral radius near $1.0$, 
resulting in painfully slow propagation of boundary information. Furthermore, it requires an 
out-of-place memory update (a "Ping-Pong" secondary buffer) to avoid data race conditions, 
effectively doubling the memory footprint and saturating DRAM bandwidth.

\subsection{The Line-Implicit Architecture (LBL TDMA)}
To drastically reduce the iteration count and memory overhead, we pivot to a Line-By-Line (LBL) implicit sweep. By treating the horizontal ($x$) direction implicitly and the vertical ($y$) direction explicitly, we collapse the 2D field into a series of coupled 1D systems:

\begin{equation}
    -a_W T_{i-1,j} + a_P T_{i,j} - a_E T_{i+1,j} = a_N T_{i,j+1} + a_S T_{i,j-1} + b
\end{equation}

For our verified Laplace case on a uniform grid, substituting the collapsed coefficients yields:

\begin{equation}
    -T_{i-1,j} + 4T_{i,j} - T_{i+1,j} = T_{i,j-1} + T_{i,j+1}
\end{equation}

This maps perfectly to our highly optimized core TDMA solver (Equation 4.5). The coefficients become $a_W = 1$, $a_E = 1$, $a_P = 4$, with the source term $d_i = T_{i,j-1} + T_{i,j+1}$.

\textbf{HPC Ramifications:} The Line-Implicit approach propagates boundary conditions across an entire row instantaneously. Crucially, it allows for \textit{in-place} memory updates (effectively becoming a Line-Gauss-Seidel method). By eliminating the $O(N \times M)$ "Ping-Pong" vector swap, we shift the algorithm from being memory-bandwidth bound to compute-bound, a vital prerequisite for scaling on modern architectures.

\section{Algorithmic Verification and Profiling}

Rigorous Test-Driven Development (TDD) via Google Test was utilized to profile the numerical efficiency of both solvers on a $50 \times 50$ Cartesian grid, subject to a stringent $L_{\infty}$ norm tolerance of $10^{-7}$.

\begin{table}[h]
\centering
\begin{tabular}{|l|c|c|l|}
\hline
\textbf{Solver Architecture} & \textbf{Max Iterations} & \textbf{Residual at Iter 1000} & \textbf{Memory Pattern} \\ \hline
Jacobi (Point-Iterative) & $> 5,400$ & $1.40 \times 10^{-3}$ & Out-of-place (Ping-Pong) \\ \hline
LBL TDMA (Line-Implicit) & $1,560$ & $9.86 \times 10^{-6}$ & In-place (No allocation) \\ \hline
\end{tabular}
\caption{Convergence profiling of 2D thermal solvers targeting $10^{-7}$ tolerance.}
\label{tab:convergence_profile}
\end{table}

\begin{figure}[htbp]
    \centering
    \includegraphics[width=0.8\textwidth]{convergence_comparison.png}
    \caption{Solver convergence comparison: $L_{\infty}$ residual vs. iteration count for the Jacobi baseline and Line-Implicit TDMA architecture.}
    \label{fig:convergence_comparison}
\end{figure}


The empirical data demonstrates a \textbf{3.4x reduction in total iterations} when 
utilizing the Line-Implicit TDMA sweep. More significantly, at the 1000-iteration mark, 
the TDMA architecture achieves a residual two orders of magnitude lower than the Jacobi 
baseline. This validates the mathematical superiority of the implicit formulation and 
establishes a highly robust numerical core for future predictive health and advanced 
thermal simulation platforms.

\begin{figure}[htbp]
    \centering
    \includegraphics[width=\textwidth]{thermal_verification_composite.png}
    \caption{Multi-modal verification of the steady-state thermal field: (Left) 2D spatial distribution showing linear propagation; (Right) 1D center-line profile validating Dirichlet boundary condition matching.}
    \label{fig:thermal_composite}
\end{figure}
